\textbf{P2)} Un deslizador de masa $0.175\ \mathrm{kg}$ se mueve sobre una pista de aire horizontal sin fricción, unido a un resorte ideal de constante elástica $k = 155\ \mathrm{N/m}$.

En cierto instante, el deslizador tiene una rapidez de $0.815\ \mathrm{m/s}$ y se encuentra a $3.00\ \mathrm{cm}$ de su posición de equilibrio.

\begin{enumerate}[a)]
  \item Calcule la amplitud $A$ del movimiento armónico simple.
  \item Determine la frecuencia angular $\omega$ y el período $T$.
  \item Escriba las ecuaciones $x(t)$, $v(t)$ y $a(t)$ suponiendo que en $t = 0$ el deslizador está en su posición de equilibrio moviéndose en sentido positivo.
\end{enumerate}

\vspace{0.3cm}